% !PRINTED 239
\weberpage{239}
But we have
\[
\varphi_2^{-1} \varphi_1 \varphi_2 = \begin{pmatrix} 0, & e^{\frac{\pi i}{4}} \\ e^{-\frac{\pi i}{4}}, & 0 \end{pmatrix}
\]
\[
\varphi_1^{-1} \varphi_2 \varphi_1 = \begin{pmatrix} 0, & e^{\frac{\pi i}{4}} \\ e^{-\frac{\pi i}{4}}, & 0 \end{pmatrix},
\]
\[
\begin{aligned}
&c\varphi_1 = \varphi_1 d,\quad d\varphi_1 = \varphi_1 c,\quad cd\varphi_1 = \varphi_1 dc\\
&c\varphi_2 = \varphi_2 dc,\quad d\varphi_2 = \varphi_2 d,\quad cd\varphi_2 = \varphi_2 c,
\end{aligned}
\]
and from these formulae it follows that the groups (11) are transformed into $D_2'$ by transformation with $\varphi_1$ and $\varphi_2$.

In the determination of the remaining polyhedral groups of the second kind the following general remarks are of use.
The inversion $j = \begin{pmatrix} i, & 0 \\ 0, & i \end{pmatrix}$ is a similarity substitution (\S.\,37) and therefore commutes with every other substitution, and consequently from every polyhedral group of the first kind we can derive at least one group of the second kind:
\weq{12}{P + Pj.}
In order to find the other groups of this kind that may still exist, we recall that by \gesperrt{Sylow}'s theorem (\S.\,29, I.) every group $G$ must contain a group whose degree is the highest power of 2 that divides the degree of $G$. But the extended tetrahedral, octahedral and icosahedral groups have the degrees 24, 48, 120, and must therefore contain a divisor of degree 8, 16, 8.

Now let $T$ be the tetrahedral, $O$ the octahedral and $J$ the icosahedral group. The groups $T$ and $J$ contain a four-group $D_2$ but no substitution of the fourth order, while $O$ contains a dihedral group $D_4$ but no substitution of the eighth order. Each extended polyhedral group must therefore contain one of the extended dihedral groups considered under III.
% [sic] the German prints \"o here (Di\"odergruppe) vs \"e elsewhere on this page (Di\"edergruppen) — reproduced as read in the German column.

Of the extended dihedral groups, however, $D_2''$ arises from $D_2$ by inversion, while
\[
\begin{aligned}
D_2' = &\begin{pmatrix}1, & 0\\ 0, & 1\end{pmatrix},\ \begin{pmatrix}\sqrt{i}, & 0\\ 0, & +i\sqrt{i}\end{pmatrix},\ \begin{pmatrix}i, & 0\\ 0, & -i\end{pmatrix},\ \begin{pmatrix}i\sqrt{i}, & 0\\ 0, & +\sqrt{i}\end{pmatrix},\\
&\begin{pmatrix}0, & -1\\ 1, & 0\end{pmatrix},\ \begin{pmatrix}0, & -i\sqrt{i}\\ \sqrt{i}, & 0\end{pmatrix},\ \begin{pmatrix}0, & i\\ i, & 0\end{pmatrix},\ \begin{pmatrix}0, & -\sqrt{i}\\ i\sqrt{i}, & 0\end{pmatrix}
\end{aligned}
\]
%CONT>>
